\section{Related Work}
\label{sec:related}

\paragraph{LLM offensive capability against sandboxes.}
SandboxEscapeBench~\cite{sandboxescapebench} evaluates whether frontier models can
break out of containerized environments spanning orchestration, container-runtime,
and kernel-level vulnerabilities, reporting that frontier models reliably escape common
misconfigurations such as privileged containers, exposed Docker sockets, and writable host mounts. Our work reuses these
scenarios but inverts the question, measuring the effect of defenses rather than
the raw offensive ceiling.

\paragraph{Defenses as a capability ceiling.}
PACEbench~\cite{pacebench} evaluates LLM agents against web applications protected by
production WAFs (ModSecurity CRS, Naxsi, Coraza) and reports that none of seven frontier
models bypassed them. Whether the same ``ceiling effect'' holds for
container/sandbox escape is untested; we evaluate this hypothesis directly.

\paragraph{Container hardening guidance.}
Operational guidance for container hardening (capability dropping, seccomp
profiles, read-only root filesystems, user-namespace remapping, minimal images,
and userspace kernels such as gVisor) is codified in community benchmarks and
vendor documentation~\cite{cisdocker,gvisor}. These sources prescribe controls but
do not quantify their marginal effect against an autonomous LLM attacker, which is
the gap we address.

\paragraph{Evaluation tooling.}
We build on the Inspect AI evaluation framework~\cite{inspectai}, extending it with
a per-control hardening wrapper and a reference-exploit verification harness.
